% ---------------------------------------------------------------------------
% Skyren — mmWave Radar Study
% Build:  lualatex mmwave_radar_study.tex   (dvakrát)
% ---------------------------------------------------------------------------
\documentclass[czech]{skyren}
\usepackage{amsmath}
\usepackage{amssymb}   % \varnothing pro průměr

\skyrendoctitle{mmWave Radar Study}
\skyrenstudytype{TECHNICKÁ STUDIE}
\skyrentitle{Levný radarový dálkoměr pro rychlé cíle}
\skyrensubtitle{Od rozebraného laserového dálkoměru k~24GHz FMCW senzoru
o~průměru 35\,mm za jednotky eur}

\skyrenaudience{Vývojář elektroniky nebo konstruktér se základy
vysokofrekvenční techniky (návrh patch antén, vedení, polarizace), který
radar dosud nenavrhoval. Tři teoretické díly (signál, odraz, anténa)
vysvětlují látku od základů i~s~odvozením; kdo ji zná, může číst jen souhrny
na začátku a~na konci dílů a~pokračovat dílem Návrh.}

\skyrenmeta{%
  \metaitem{Projekt}{mmWave Radar}\hfill
  \metaitem{Revize}{R2}\hfill
  \metaitem{Datum}{Září 2026}\hfill
  \metaitem[0.2]{Klasifikace}{Interní}%
}

\begin{document}

\skyrenmaketitle

\begin{skyrensummary}
Vzdálenost k~plošnému cíli do 5\,m při vzájemné rychlosti až 240\,m/s
změří \textbf{24GHz FMCW radar s~trojúhelníkovou rampou}, jednou anténou
o~průměru 35\,mm a~tištěnou čočkou; materiál vychází na \textbf{zhruba
4--7\,€}. Na hladkých površích (plech, hladký beton, hoblované dřevo) je
chyba pod 1\,cm a~radar měří kolmou vzdálenost k~rovině. Na drsných površích
(hrubý beton, cihla) rozhoduje \textbf{speckle}: po průměrování přes 1,2\,m
přiblížení vychází chyba kolem 3\,cm u~kolmé zdi, ale 5--6\,cm při náklonu
15$^\circ$. Tam zadání 5\,cm splní až užší svazek. Otevřená je dokumentace
a~cena čipu; chceme je získat rozebráním modulu HLK-LD2450.
\end{skyrensummary}

\begin{skyrennote}[Jak číst tento dokument]
\textbf{Boxy se štítkem} označují typ sdělení: \textsc{poznámka} kontext
navíc, \textsc{úskalí} častou chybu, \textsc{příklad} konkrétní výpočet.
\textsc{Na jeden pohled} otevírá díl, \textsc{co z~toho plyne} ho uzavírá.
\skyrencmd{Mono písmo} značí hodnotu nebo označení dílu. Revize R2 nahrazuje
R1: zadání se mezitím zpřesnilo a~většina závěrů R1 už neplatí. Čísla v~R2
pocházejí z~výpočtů a~simulací, ne z~měření.
\end{skyrennote}

\tableofcontents
\vspace{10mm}

% ===========================================================================
\skyrenchapter[\skyrenmodeexplanation]{Východiska}{Zadání a proč ne optika}
% ===========================================================================

\begin{skyrenglance}
Rozebraný Parkside PLEM 20 A4 je fázový měřák s~uzavřeným ASICem --- slepá
ulička. Zadání 240\,m/s pak vyřadilo optiku jako princip: limitem přestává
být přesnost jednoho měření a~stává se jím vzorkovací frekvence. Rozhodující
pro návrh jsou ale až další dvě podmínky --- cena a~velikost.
\end{skyrenglance}

\section{Co jsme rozebrali}
\skyrenlead{Teardown laserového dálkoměru Parkside a~důvod, proč se v~něm
nedá pokračovat.}

Výchozím bodem byl rozebraný dálkoměr \skyrencmd{Parkside PLEM 20 A4}
(model \skyrencmd{HG06724}, výroba 2020/06). Měřicí modul je samostatná deska
zhruba 13\,$\times$\,25\,mm: laserová dioda v~pouzdře TO, přijímací lavinová
fotodioda ve~stíněné komoře, krystal a~jeden ASIC v~pouzdře QFN s~potiskem
\skyrencmd{LGSM4 / N1PA45.00-24}. K~hlavní desce vede osmipinový plochý spoj.

Přístroj neměří dobu letu pulzu, ale \skyrenterm{fázový posuv} spojitě
modulovaného svazku. Celý měřicí řetězec je uvnitř ASICu se zákaznickým
označením bez veřejného datového listu. Odposlech plochého spoje by dal
kontrolu nad rozhraním, ne nad měřením: frekvence i~integrační čas zůstávají
zadrátované v~čipu.

\section{Zadání}
\skyrenlead{Pět požadavků. Poslední dva, cena a~velikost, rozhodly o~návrhu
nejvíc.}

Zadání se během studie zpřesňovalo. Tabulka uvádí jeho konečný stav:

\begin{skyrentable}{3}{Konečné zadání. Cena a~velikost jsou priority, které
určil zadavatel.}
\skyrenth{Parametr} & \skyrenth{Hodnota} & \skyrenth{Poznámka} \\
\midrule
Cíl & jeden, plošný & kov, plech, beton, zeď, dřevo \\
Přesnost & $\approx$\,5\,cm na 5\,m & rozšíření: $<$\,3\,cm do 30\,m \\
Rychlost & až 240\,m/s & vzájemná, cíl se přibližuje \\
Cena & ideálně 25\,€, max. 50\,€ & materiál na kus; \textbf{priorita č.~1} \\
Velikost & kruh $\varnothing$\,35\,mm, co nejmenší & \textbf{priorita č.~2} \\
\end{skyrentable}

Důležité je, že jde o~\emph{jeden} cíl. Nepotřebujeme oddělit dva odrazy
blízko sebe (\skyrenterm{rozlišení}), stačí přesně určit polohu jednoho
(\skyrenterm{přesnost}). Teoretické díly ukazují, že přesnost může být
o~řád lepší než rozlišení.

\section{Co udělá rychlost 240\,m/s}
\skyrenlead{Převod rychlosti na požadavek na vzorkovací frekvenci.}

Mezi dvěma vzorky urazí cíl dráhu, kterou žádné přesnější měření
neodstraní --- jen vyšší frekvence:

\begin{skyrentable}{3}{Dráha mezi dvěma vzorky při 240\,m/s.}
\skyrenth{Vzorkovací frekvence} & \skyrenth{Perioda} & \skyrenth{Dráha mezi vzorky} \\
\midrule
100\,Hz & 10\,ms & 2,40\,m \\
1\,kHz & 1\,ms & 24\,cm \\
4\,kHz & 250\,\textmu s & 6\,cm \\
10\,kHz & 100\,\textmu s & 2,4\,cm \\
\end{skyrentable}

Z~toho plynou tři požadavky: vzorkovací frekvence aspoň 1\,kHz,
integrační čas jednoho měření pod 1\,ms a~latence celé cesty v~jednotkách
milisekund. Měření staré 5\,ms je 1,2\,m mimo.

\begin{skyrengotcha}[Přesnost z~datového listu nestačí]
Údaj „$\pm$\,3\,cm“ v~datovém listu popisuje rozptyl jednoho měření za
ideálních podmínek. Při 240\,m/s bývá větší chybou stáří vzorku. Tu datový
list neuvede, protože závisí na vaší rychlosti --- spočítejte ji z~frekvence
a~latence.
\end{skyrengotcha}

\section{Proč ne optika}
\skyrenlead{Jediný vyhovující optický dálkoměr je o~řád dražší, než
dovoluje rozpočet.}

Přímé měření doby letu (dToF) splní požadavek na frekvenci jen u~výjimek.
\skyrencmd{LightWare SF30/D} měří katalogově 0,2 až 200\,m s~frekvencí až
20\,000 měření za sekundu. Rychlostí by zadání splnil, cenou ale patří mezi
měřicí techniku, ne mezi součástky. Levné senzory (\skyrencmd{VL53Lxx},
\skyrencmd{TMF8828}, \skyrencmd{TF-Luna}) běží na desítkách až stovkách Hz.
\skyrencmd{Broadcom AFBR-S50MV85G} zvládne až 3000 snímků za sekundu, ale
dosah má kolem 10\,m. S~rozpočtem 25\,€ proto zbývá radar.

\begin{skyrentakeaway}
Optika na tomhle zadání neselhala na přesnosti, ale na kombinaci frekvence
a~ceny. O~dalším návrhu rozhodly cena a~velikost, ne fyzikální meze měření.
\end{skyrentakeaway}

% ===========================================================================
\skyrenchapter[\skyrenmodeexplanation]{Teorie: signál}{Signál FMCW}
% ===========================================================================

\begin{skyrenglance}
FMCW převádí vzdálenost na kmitočet. Vzdálenost vypadne z~polohy vrcholu
ve spektru jedné rampy; rozlišení určuje šířka pásma, přesnost odstup
signálu od šumu. Pohyb cíle posune odhad o~desítky centimetrů, trojúhelníková
rampa ten posuv vyruší. U~blízkých cílů rozhoduje vazba vysílače do
přijímače a~blikavý šum.
\end{skyrenglance}

\section{Chirp a~rozdílová frekvence}
\skyrenlead{Model signálu, ze kterého plyne všechno ostatní.}

\skyrenterm{FMCW} (frequency-modulated continuous wave) vysílá spojitý signál,
jehož kmitočet lineárně roste; jedné takové rampě se říká \skyrenterm{chirp}.
Se sklonem $S = B/T$, kde $B$ je šířka přeladění a~$T$ délka chirpu, má
vysílaný signál fázi
\[
\varphi_\mathrm{t}(t) = 2\pi\left(f_0 t + \tfrac{1}{2} S t^2\right)
\]
Odraz od cíle ve vzdálenosti $R$ přijde zpožděný o~dobu letu $\tau = 2R/c$.
Směšovač vynásobí vyslaný a~přijatý signál a~dolní propust nechá jen
rozdílovou složku:
\[
s_\mathrm{b}(t) = A\cos\!\left(2\pi S\tau\,t + 2\pi f_0\tau - \pi S\tau^2\right)
\]
Je to harmonický signál s~\skyrenterm{rozdílovou frekvencí}
\[
f_\mathrm{b} = S\tau = \frac{2RB}{cT}
\]
a~s~fází $4\pi R/\lambda$, která se mění už při posunu cíle o~zlomek vlnové
délky. Člen $\pi S\tau^2$ je na 5\,m pod 0,1\,rad a~konstantní, takže ho lze
zanedbat.

\begin{skyrenfigure}{Princip FMCW. Vyslaný chirp (plná čára) stoupá lineárně;
odraz (čárkovaně) je jeho kopie posunutá o~dobu letu $\tau$. Svislá vzdálenost
mezi čarami je rozdílová frekvence $f_\mathrm{b}$ --- po celou rampu konstantní
a~přímo úměrná vzdálenosti cíle.}{fig:fmcw}
\begin{tikzpicture}[x=1.2mm, y=1.2mm, line width=1pt,
  every node/.style={font=\fontsize{8}{10}\selectfont}]
  \draw[skyrenNeutral600, line width=0.6pt, ->] (0,0) -- (78,0)
    node[anchor=north east] {čas};
  \draw[skyrenNeutral600, line width=0.6pt, ->] (0,0) -- (0,32)
    node[anchor=south west] {kmitočet};
  \draw[skyrenInk] (0,2) -- (30,26);
  \draw[skyrenInk] (30,2) -- (60,26);
  \draw[skyrenNeutral600, dashed, line width=0.9pt] (7,2) -- (37,26);
  \draw[skyrenNeutral600, dashed, line width=0.9pt] (37,2) -- (60,20.4);
  \draw[skyrenRed, line width=1pt] (22,13.6) -- (22,19.6);
  \node[anchor=west, skyrenRed] at (23.5,16.6) {$f_\mathrm{b}$};
  \draw[skyrenNeutral600, line width=0.6pt] (0,-4) -- (7,-4);
  \draw[skyrenNeutral600, line width=0.6pt] (0,-2.5) -- (0,-5.5);
  \draw[skyrenNeutral600, line width=0.6pt] (7,-2.5) -- (7,-5.5);
  \node[anchor=north] at (3.5,-5) {$\tau$};
  \draw[skyrenNeutral600, line width=0.6pt] (66,2) -- (66,26);
  \draw[skyrenNeutral600, line width=0.6pt] (64.5,2) -- (67.5,2);
  \draw[skyrenNeutral600, line width=0.6pt] (64.5,26) -- (67.5,26);
  \node[anchor=west] at (68,14) {$B$};
  \draw[skyrenInk] (0,-13) -- (7,-13);
  \node[anchor=west] at (8.5,-13) {vyslaný chirp};
  \draw[skyrenNeutral600, dashed, line width=0.9pt] (34,-13) -- (41,-13);
  \node[anchor=west] at (42.5,-13) {přijatý odraz};
\end{tikzpicture}
\end{skyrenfigure}

Na obr.~\ref{fig:fmcw} je vidět, že rozdíl mezi čarami je po celou rampu
konstantní. Pro návrhový profil ($B = 250$\,MHz, $T = 10$\,\textmu s) vychází
$f_\mathrm{b}$ na 1\,m 167\,kHz, na 5\,m 833\,kHz a~na 8\,m 1,33\,MHz.

\section{Od vzorků ke vzdálenosti}
\skyrenlead{Okno, FFT, doplnění nulami a~interpolace vrcholu --- a~kolik
která část přispívá k~chybě.}

Mezifrekvenční signál se během chirpu navzorkuje; při 4\,MSPS a~10\,\textmu s
to je $N = 40$ vzorků. Vzorkování \skyrenterm{I/Q} (dva kanály posunuté
o~90$^\circ$) dává komplexní signál, který rozliší znaménko kmitočtu. To je
potřeba pod 0,23\,m, kde dopplerovský posuv převýší rozdílovou frekvenci
(viz dále).

Vzorky se vynásobí \skyrenterm{oknem}. Hannovo okno potlačí postranní laloky
spektra na $-31$\,dB, takže silná vazba vysílače do přijímače nebo druhý
odraz neprosáknou do buňky cíle. Cenou je dvakrát širší hlavní lalok
a~šumová šířka 1,5 buňky místo jedné.

Pak následuje rychlá Fourierova transformace (FFT). Doplnění nulami ze 40 na
64 nebo 128 bodů zhustí vzorkování spektra, rozlišení ale nezlepší.
Vrchol leží mezi body FFT; jeho polohu dopočítá
\skyrenterm{parabolická interpolace} logaritmu amplitudy ve třech bodech
$k-1, k, k+1$ s~hodnotami $\alpha, \beta, \gamma$ v~dB:
\[
\delta = \frac{1}{2}\cdot\frac{\alpha-\gamma}{\alpha-2\beta+\gamma}, \qquad
\hat R = (k+\delta)\cdot\frac{c\,f_\mathrm{s}}{2\,S\,M}
\]
kde $f_\mathrm{s}$ je vzorkovací frekvence a~$M$ délka FFT. Simulace
s~Hannovým oknem dala systematickou chybu interpolace nejvýš 2\,mm pro
$M = 64$ a~0,2\,mm pro $M = 128$. Interpolace tedy do rozpočtu chyby
nepřispívá.

\section{Rozlišení a~přesnost}
\skyrenlead{Rozlišení určuje šířka pásma, přesnost na jeden cíl určuje odstup
signálu od šumu.}

\skyrenterm{Rozlišení} je nejmenší vzdálenost, na které radar oddělí dva
cíle. Závisí jen na šířce přeladění:
\[
\Delta R = \frac{c}{2B}
\]
Ve volném pásmu \skyrencmd{24,00--24,25\,GHz} je $B = 250$\,MHz, takže
$\Delta R = 60$\,cm. Lepší rozlišení by vyžadovalo jiné pásmo: 4\,GHz
v~pásmu 60\,GHz dává 3,75\,cm.

Pro jeden cíl je podstatná \skyrenterm{přesnost}. Teoretickou nejmenší
chybu udává \skyrenterm{Cramér--Raova mez} pro odhad kmitočtu harmonického
signálu v~šumu:
\[
\sigma_R \geq \frac{\sqrt{6}}{2\pi}\cdot\frac{\Delta R}{\sqrt{\mathrm{SNR}}}
\approx 0{,}39\,\frac{\Delta R}{\sqrt{\mathrm{SNR}}}
\]
kde SNR je odstup signálu od šumu po FFT (energie signálu za chirp
k~spektrální hustotě šumu). Hannovo okno s~parabolickou interpolací dosahuje
podle simulace $\sigma_R \approx 0{,}6\,\Delta R/\sqrt{\mathrm{SNR}}$.

\begin{skyrenexample}
Buňka 60\,cm, SNR 30\,dB (1000$\times$): mez dává 0,74\,cm, praktický
odhad 1,1\,cm, simulace 1,2\,cm. Při 40\,dB vychází prakticky 0,35\,cm.
Rozlišení 60\,cm tedy přesnosti na centimetry nebrání, pokud je signálu
dost a~cíl je jeden.
\end{skyrenexample}

Fáze vrcholu $4\pi R/\lambda$ nese ještě jemnější informaci, ale
s~nejednoznačností $\lambda/2$, na 24\,GHz 6,2\,mm. Hrubý odhad z~FFT není
dost přesný, aby ji rozbalil absolutně. Fáze proto dává jen
\emph{relativní} změnu vzdálenosti při spojitém sledování cíle.

\section{Dopplerův posuv a~trojúhelníková rampa}
\skyrenlead{Při 240\,m/s posune pohyb cíle odhad o~23\,cm. Střídání rampy
nahoru a~dolů ten posuv vyruší a~navíc změří rychlost.}

Přibližující se cíl zvýší kmitočet odrazu o~\skyrenterm{dopplerovský posuv}
$f_\mathrm{d} = 2v/\lambda$; při 240\,m/s a~$\lambda = 12{,}4$\,mm to je
38,6\,kHz. U~stoupající rampy se tím rozdílová frekvence zmenší,
u~klesající zvětší:
\[
f_{\mathrm{b}\uparrow} = S\tau - f_\mathrm{d}, \qquad
f_{\mathrm{b}\downarrow} = S\tau + f_\mathrm{d}
\]
Převedeno na vzdálenost je každý z~obou odhadů posunutý o
\[
\Delta R_\mathrm{D} = \frac{v\,c\,T}{\lambda\,B}
\]
což je při chirpu 10\,\textmu s \textbf{23\,cm}, při 25\,\textmu s už 58\,cm.

\begin{skyrenfigure}{Trojúhelníková rampa a~přibližující se cíl. Odraz
(čárkovaně) je zpožděný o~$\tau$ a~posunutý nahoru o~$f_\mathrm{d}$. Na
stoupající rampě se rozdílová frekvence zmenší, na klesající zvětší; průměr
obou je bez Dopplerova posuvu.}{fig:troj}
\begin{tikzpicture}[x=1.2mm, y=1.2mm, line width=1pt,
  every node/.style={font=\fontsize{8}{10}\selectfont}]
  \draw[skyrenNeutral600, line width=0.6pt, ->] (0,0) -- (70,0)
    node[anchor=north east] {čas};
  \draw[skyrenNeutral600, line width=0.6pt, ->] (0,0) -- (0,33)
    node[anchor=south west] {kmitočet};
  \draw[skyrenInk] (0,2) -- (30,26) -- (60,2);
  \draw[skyrenNeutral600, dashed, line width=0.9pt] (6,4.5) -- (36,28.5) -- (60,9.3);
  \draw[skyrenRed, line width=1pt] (20,15.7) -- (20,18);
  \node[anchor=east, skyrenRed] at (19,16.9) {$f_{\mathrm{b}\uparrow}$};
  \draw[skyrenRed, line width=1pt] (48,11.6) -- (48,18.9);
  \node[anchor=west, skyrenRed] at (49,15.2) {$f_{\mathrm{b}\downarrow}$};
  \draw[skyrenInk] (0,-6) -- (7,-6);
  \node[anchor=west] at (8.5,-6) {vyslaná rampa};
  \draw[skyrenNeutral600, dashed, line width=0.9pt] (34,-6) -- (41,-6);
  \node[anchor=west] at (42.5,-6) {odraz od přibližujícího se cíle};
\end{tikzpicture}
\end{skyrenfigure}

Z~dvojice odhadů $R_\uparrow = R - \Delta R_\mathrm{D}$
a~$R_\downarrow = R + \Delta R_\mathrm{D}$ (obr.~\ref{fig:troj}) plyne
\[
R = \frac{R_\uparrow + R_\downarrow}{2}, \qquad
v = \frac{\lambda B}{2cT}\,\bigl(R_\downarrow - R_\uparrow\bigr)
\]
Vzdálenost platí pro střed dvojice chirpů, tam patří i~časová značka. Během
jednoho chirpu se cíl posune o~$vT = 2{,}4$\,mm, což je zanedbatelné.
Dopplerovská FFT přes více chirpů, obvyklá u~automobilových radarů, tu
nefunguje: s~periodou chirpů kolem 25\,\textmu s je nejvyšší jednoznačná
rychlost $\lambda/4T_\mathrm{c} \approx 124$\,m/s a~240\,m/s by se přetočila.

\section{Šum a~odstup signálu od šumu}
\skyrenlead{Tepelný šum v~šířce jedné buňky a~proč fázový šum oscilátoru na
metrech skoro nevadí.}

Šum připadající na jednu buňku spektra je tepelný šum v~šumové šířce okna,
zvětšený o~šumové číslo přijímače $F$:
\[
N = kT_0 \cdot \frac{\mathrm{ENBW}}{T} \cdot F
\quad\Rightarrow\quad
N_\mathrm{dBm} = -174 + 10\log_{10}\frac{1{,}5}{T} + F_\mathrm{dB}
\]
Pro chirp 10\,\textmu s a~šumové číslo 15\,dB (odhad, datový list čipu
nemáme) vychází \mbox{$N \approx -107$\,dBm}.

Fázový šum oscilátoru se ve směšovači z~velké části odečte, protože vysílaný
i~referenční signál pocházejí ze stejného oscilátoru. Zbytek je úměrný
$2\sin(\pi f_\mathrm{m}\tau)$, kde $f_\mathrm{m}$ je odstup od nosné. Na
5\,m je $\tau = 33$\,ns a~potlačení vychází $-14$\,dB při 1\,MHz a~$-34$\,dB
při 100\,kHz. Na krátké vzdálenosti tedy kvalita oscilátoru rozhoduje málo.

\section{Blízký dosah: vazba, blikavý šum a~horní propust}
\skyrenlead{Silná vazba vysílače do přijímače, blikavý šum a~velká dynamika
se řeší jedním filtrem --- a~ten má mít pro plošný cíl sklon 20\,dB na
dekádu, ne 40.}

S~jednou anténou pro vysílání i~příjem teče část vysílaného výkonu přímo do
přijímače: přes konečnou izolaci vazebníku (typicky kolem 25\,dB), přes odraz
na anténě a~přes odrazy na čočce. Tato vazba sedí u~nulové rozdílové
frekvence a~může dosáhnout $-15$ až $-20$\,dBm na vstupu přijímače.
Kompresní bod přijímače integrovaného čipu neznáme; pokud je kolem
$-15$\,dBm, hrozí zahlcení. Proto je u~jedné antény klíčové dobré
impedanční přizpůsobení ($S_{11} < -20$\,dB).

Blikavý šum (1/$f$) směšovače roste pod zhruba 100\,kHz. Krátký chirp drží
rozdílovou frekvenci i~na 1\,m (167\,kHz) nad touto oblastí.

Výkon odrazu od plošného cíle klesá s~$R^2$ (díl Odraz od plošného cíle), amplituda tedy
s~$1/R$, a~protože $f_\mathrm{b} \propto R$, klesá amplituda úměrně
$1/f_\mathrm{b}$. Horní propust prvního řádu s~mezním kmitočtem nad pásmem
měření má v~pásmu zisk úměrný $f$, tedy \textbf{+20\,dB na dekádu}: srovná
úroveň signálu přes celý rozsah vzdáleností a~současně potlačí vazbu
u~nulového kmitočtu. Pro bodový cíl ($R^{-4}$) by bylo potřeba +40\,dB na
dekádu.

\begin{skyrentakeaway}
Vzdálenost dá jedna rampa a~jedna malá FFT; interpolace vrcholu je přesná na
milimetry. Dopplerův posuv 23\,cm vyruší trojúhelníková rampa. Přesnost pak
určuje odstup signálu od šumu --- a~u~plošného cíle, jak ukáže další díl,
hlavně geometrie odrazu.
\end{skyrentakeaway}

% ===========================================================================
\skyrenchapter[\skyrenmodeexplanation]{Teorie: odraz}{Odraz od plošného cíle}
% ===========================================================================

\begin{skyrenglance}
Hladká plocha se chová jako zrcadlo: odraz přichází z~malé skvrny kolem
paty kolmice a~radar změří kolmou vzdálenost na milimetry. Drsná plocha
rozptyluje: přispívá celá stopa svazku, odhad se posune a~náhodně kolísá
(speckle). Rozhoduje drsnost vůči vlnové délce, ne materiál.
\end{skyrenglance}

\section{Tři modely odrazu}
\skyrenlead{Bodový cíl, zrcadlová plocha a~drsná plocha mají tři různé
radarové rovnice.}

\textbf{Bodový cíl} s~radarovým průřezem $\sigma$ je výchozí případ radarové
rovnice:
\[
P_\mathrm{r} = \frac{P_\mathrm{t}\,G^2\lambda^2\,\sigma}{(4\pi)^3 R^4}
\]

\textbf{Hladká rovina} větší než první Fresnelova zóna se chová jako zrcadlo.
Radar v~ní vidí svůj obraz ve vzdálenosti $2R$, takže
\[
P_\mathrm{r} = \frac{P_\mathrm{t}\,G^2\lambda^2\,|\Gamma|^2}{(4\pi)^2 (2R)^2}
\]
kde $\Gamma$ je činitel odrazu materiálu. Odraz vzniká v~\skyrenterm{první
Fresnelově zóně}, skvrně o~poloměru $\sqrt{\lambda R/2}$ kolem paty kolmice
z~radaru na rovinu. Na 5\,m to je 18\,cm, zatímco stopa svazku 22$^\circ$ má
poloměr kolem 1\,m. Zeď je tedy pro radar „nekonečné zrcadlo“.

\textbf{Drsná rovina} rozptyluje difúzně a~každý kousek plochy přispívá
podle \skyrenterm{koeficientu zpětného rozptylu} $\sigma^0$ (m$^2$ na
m$^2$ plochy). Efektivní průřez je $\sigma = \sigma^0 R^2 \Omega_\mathrm{ef}$,
kde $\Omega_\mathrm{ef}$ je prostorový úhel svazku vážený dvoucestným
diagramem; pro gaussovský svazek s~šířkou $\theta_3$ platí
$\Omega_\mathrm{ef} = \pi\theta_3^2/(8\ln 2)$. Protože zisk antény
$G \propto 1/\theta_3^2$, vychází
\[
P_\mathrm{r} \propto \frac{G}{R^2}
\]
Oba plošné případy tedy klesají s~$R^2$, ne s~$R^4$. U~drsné plochy navíc
větší anténa signál skoro nezlepší, protože užší svazek ozáří menší plochu.

\section{Hladký, nebo drsný?}
\skyrenlead{O~chování plochy rozhoduje drsnost ve srovnání s~vlnovou délkou.
Na 24\,GHz leží hranice kolem 1,5\,mm.}

Nerovnosti o~střední výšce $h$ rozfázují odražené vlny. Zrcadlová
(koherentní) složka odrazu klesne podle \skyrenterm{Rayleighova kritéria}
o~faktor
\[
\rho = \exp\!\left[-\left(\frac{4\pi h\cos\theta_\mathrm{i}}{\lambda}\right)^2\right]
\]
kde $\theta_\mathrm{i}$ je úhel dopadu. Na 24\,GHz to dává pokles o~1,1\,dB
při $h = 0{,}5$\,mm, 4,4\,dB při 1\,mm, 18\,dB při 2\,mm a~40\,dB při 3\,mm.
Obvyklá hranice je $h \approx \lambda/8 = 1{,}6$\,mm.

\begin{skyrentable}{4}{Typické povrchy na 24\,GHz. Drsnosti a~permitivity
jsou orientační hodnoty, ne měření.}
\skyrenth{Povrch} & \skyrenth{Drsnost $h$} & \skyrenth{Chování} & \skyrenth{$|\Gamma|$ kolmo} \\
\midrule
plech & $\ll$\,0,1\,mm & zrcadlové & 1 (0\,dB) \\
hladký beton, omítka & $\approx$\,0,5--1\,mm & převážně zrcadlové & $\approx$\,0,40 ($-8$\,dB) \\
hoblované dřevo & $<$\,0,5\,mm & zrcadlové & $\approx$\,0,17 ($-15$\,dB) \\
sádrokarton & $<$\,0,5\,mm & zrcadlové & $\approx$\,0,25 ($-12$\,dB) \\
řezané dřevo & $\approx$\,1--2\,mm & smíšené & $\approx$\,0,17 \\
hrubý beton, cihla, kámen & $\approx$\,2--5\,mm & difúzní & --- \\
\end{skyrentable}

Činitel odrazu při kolmém dopadu je $\Gamma = (1-\sqrt{\varepsilon_\mathrm{r}})/(1+\sqrt{\varepsilon_\mathrm{r}})$.
Pro beton s~$\varepsilon_\mathrm{r} \approx 5{,}5$ vychází 0,40, pro suché
dřevo s~$\varepsilon_\mathrm{r} \approx 2$ jen 0,17. Dřevo je tedy nejslabší
zrcadlo, ale rozpočet výkonu v~dílu Návrh ukáže, že i~tak zbývá přes
40\,dB odstupu od šumu.

\section{Zrcadlový odraz měří kolmou vzdálenost}
\skyrenlead{Hladká plocha nemá žádný posun od stopy svazku. Při náklonu ale
měří jinou vzdálenost, než je vzdálenost v~ose, a~slábne.}

Protože zrcadlový odraz přichází jen z~první Fresnelovy zóny kolem paty
kolmice, radar změří \emph{kolmou vzdálenost k~rovině} $R_\perp = R\cos\alpha$,
kde $\alpha$ je náklon plochy vůči ose senzoru. Stopa svazku se neuplatní
a~chyba je daná jen šumem.

Při náklonu leží pata kolmice mimo osu svazku a~odraz zeslábne podle
dvoucestného diagramu. Pro svazek 22$^\circ$ je to $-6$\,dB při
11$^\circ$, $-11$\,dB při 15$^\circ$, $-20$\,dB při 20$^\circ$ a~$-31$\,dB
při 25$^\circ$. Hladký plech nakloněný o~víc než zhruba 20$^\circ$ proto
nevrátí skoro nic --- nemá difúzní složku, která by zbyla.

\begin{skyrennote}[Kterou vzdálenost vlastně měříme]
Při náklonu 15$^\circ$ je na 5\,m rozdíl mezi kolmou vzdáleností a~vzdáleností
v~ose kolem 17\,cm. Není to chyba měření. U~hladkých ploch radar měří
kolmou vzdálenost, u~drsných zhruba vzdálenost v~ose. Aplikace musí vědět,
kterou z~nich potřebuje.
\end{skyrennote}

\section{Stopa svazku a~posun odhadu}
\skyrenlead{U~drsné plochy přispívá celá stopa. Její zakřivení vůči radaru
posune odhad za skutečnou vzdálenost.}

\begin{skyrenfigure}{Stopa svazku na zdi (bokorys, úhel zveličený). Kružnice
konstantní vzdálenosti (červeně) protíná zeď jen v~ose; okraje stopy leží
dál o~$\Delta R$. U~drsné plochy padnou všechny příspěvky do jedné buňky
a~posunou odhad za kolmou vzdálenost.}{fig:stopa}
\begin{tikzpicture}[x=1.1mm, y=1.1mm, line width=1pt,
  every node/.style={font=\fontsize{8}{10}\selectfont}]
  \draw[skyrenInk, fill=skyrenNeutral400] (-6,-3) rectangle (0,3);
  \node[anchor=east] at (-7,0) {radar};
  \draw[skyrenNeutral600, line width=0.7pt] (0,0) -- (60,21.84);
  \draw[skyrenNeutral600, line width=0.7pt] (0,0) -- (60,-21.84);
  \draw[skyrenNeutral600, dashed, line width=0.7pt] (0,0) -- (60,0);
  \node[anchor=south] at (30,0.5) {$R$};
  \draw[skyrenRed, line width=0.9pt] (56.38,-20.52) arc[start angle=-20, end angle=20, radius=60];
  \draw[skyrenInk, line width=2.2pt] (60,-27) -- (60,27);
  \node[anchor=west] at (61.5,-24) {zeď};
  \draw[skyrenInk, line width=0.7pt] (56.38,20.52) -- (60,21.84);
  \node[anchor=south east] at (57,21.5) {$\Delta R$};
  \draw[skyrenNeutral600, line width=0.6pt] (14,0) arc[start angle=0, end angle=20, radius=14];
  \node[anchor=west] at (15,2.6) {$\theta/2$};
\end{tikzpicture}
\end{skyrenfigure}

Vzdálenost k~bodu plochy ve směru $(\theta_x, \theta_y)$ od osy je pro plochu
nakloněnou o~$\alpha$ kolem osy $y$ přibližně
\[
R(\theta_x,\theta_y) \approx R\left[1 + \theta_x\tan\alpha
  + \theta_x^2\left(\tfrac{1}{2}+\tan^2\alpha\right) + \tfrac{1}{2}\theta_y^2\right]
\]
Rozdíly vzdáleností ve stopě jsou mnohem menší než buňka 60\,cm, takže vrchol
FFT leží ve výkonově váženém průměru vzdáleností. Pro gaussovský dvoucestný
diagram s~rozptylem $\sigma_\theta^2 = \theta_3^2/(16\ln 2)$ v~každé ose
lineární člen vypadne a~zbude \skyrenterm{posun od stopy svazku}
\[
b = R\,\sigma_\theta^2\,(1+\tan^2\alpha) = \frac{R\,\theta_3^2}{16\ln 2}\,(1+\tan^2\alpha)
\]
(obr.~\ref{fig:stopa}). U~kolmé zdi se simulace od vzorce liší o~méně než
desetinu, u~nakloněné do pětiny, zčásti kvůli statistickému rozptylu
simulace.

\begin{skyrentable}{3}{Posun odhadu od stopy svazku pro drsnou zeď kolmo
k~ose. Hodnota je deterministická a~odečítá se kalibrací.}
\skyrenth{HPBW} & \skyrenth{Posun na 5\,m} & \skyrenth{Posun na 30\,m} \\
\midrule
42$^\circ$ & 24\,cm & 1,45\,m \\
22$^\circ$ & 6,6\,cm & 40\,cm \\
14$^\circ$ & 2,7\,cm & 16\,cm \\
10$^\circ$ & 1,4\,cm & 8,2\,cm \\
\end{skyrentable}

Posun je pro známou vzdálenost a~svazek spočitatelný, takže se odečte. Náklon
ho zvětší o~$\tan^2\alpha$, při 15$^\circ$ o~7\,\%. Zbytek po kalibraci je na
5\,m se svazkem 22$^\circ$ kolem 0,5\,cm z~náklonu, k~tomu neznámý příspěvek
toho, že $\sigma^0$ závisí na úhlu dopadu a~váží stopu nesymetricky. Počítáme
s~celkem kolem 1\,cm, ale ověřené to není.

\section{Speckle}
\skyrenlead{Příspěvky drsné plochy se sčítají s~náhodnými fázemi. Odhad
vzdálenosti proto kolísá, a~u~nakloněné plochy hodně.}

Každý kousek drsné plochy vrací vlnu s~náhodnou fází. Jejich součet má
náhodnou amplitudu i~náhodnou polohu vrcholu ve spektru; tomuto jevu se
říká \skyrenterm{speckle}. Střední hodnota odhadu je posun $b$ z~předchozí
sekce, kolem ní odhad kolísá zhruba o~efektivní rozptyl vzdáleností ve stopě
\[
s_R \approx R\,\sigma_\theta\sqrt{\sigma_\theta^2 + \tan^2\alpha}
\]
U~kolmé plochy je $s_R = b$, u~nakloněné převládne lineární člen
$R\,\sigma_\theta\tan\alpha$ a~kolísání vzroste několikanásobně.
Když amplituda odrazu náhodně klesne (hluboký útlum), poloha vrcholu ujede
výrazně víc; takové měření je lepší zahodit.

\begin{skyrentable}{4}{Kolísání jednoho měření na drsné zdi ve 5\,m. Simulace
Monte Carlo s~vyřazením 10\,\% nejslabších měření, srovnaná s~modelem.}
\skyrenth{HPBW} & \skyrenth{Náklon} & \skyrenth{Simulace} & \skyrenth{Model $s_R$} \\
\midrule
22$^\circ$ & 0$^\circ$ & 7,1\,cm & 6,7\,cm \\
22$^\circ$ & 15$^\circ$ & 15,8\,cm & 16,8\,cm \\
14$^\circ$ & 0$^\circ$ & 3,0\,cm & 2,7\,cm \\
14$^\circ$ & 15$^\circ$ & 10,6\,cm & 10,2\,cm \\
10$^\circ$ & 0$^\circ$ & 1,3\,cm & 1,4\,cm \\
10$^\circ$ & 15$^\circ$ & 7,2\,cm & 7,2\,cm \\
\end{skyrentable}

Kolísání pomůže jen průměrování přes \emph{nezávislé} realizace speckle. Při
přibližování se mění relativní fáze příspěvků a~nová realizace vznikne
zhruba po každých
\[
\delta R \approx \frac{\lambda R}{4\,s_R}
\]
přiblížení, na 5\,m u~kolmé zdi a~svazku 22$^\circ$ po 24\,cm. Průměr přes
dráhu přiblížení $L$ má kolísání
\[
\sigma_\mathrm{prům} \approx \sqrt{\frac{s_R\,\lambda R}{4L}}
\]
Pro $L = 1{,}2$\,m (5\,ms při 240\,m/s) dává model 2,9\,cm u~kolmé zdi
a~4,7\,cm při náklonu 15$^\circ$; simulace přibližování dala 2,3\,cm
a~5,8\,cm.

\begin{skyrengotcha}[Průměruje se dráha, ne čas]
Speckle se obnovuje pohybem, ne časem. Při 240\,m/s dá 5\,ms dost nezávislých
měření, při 10\,m/s by stejné okno vystačilo jen na 5\,cm dráhy a~kolísání by
zůstalo skoro celé. Když senzor stojí, je speckle zamrzlý a~chyba se stane
pevnou chybou daného místa. Pásmo 250\,MHz je navíc příliš úzké na to, aby
pomohlo střídání kmitočtu.
\end{skyrengotcha}

\section{Tenké dielektrické desky}
\skyrenlead{Dřevěná deska nebo sádrokarton odrážejí z~přední i~zadní strany
a~oba odrazy padnou do jedné buňky.}

Část vlny projde do materiálu s~nízkou permitivitou a~odrazí se od zadní
strany. Ta se radaru jeví posunutá o~$\sqrt{\varepsilon_\mathrm{r}}\,d$; pro
2cm překližku s~$\varepsilon_\mathrm{r} \approx 2$ o~2,8\,cm. Oba odrazy
leží v~jedné buňce a~interferují podle tloušťky desky, takže vrchol se posune
o~část této vzdálenosti. Útlum ve vlhkém dřevě zadní odraz zeslabí. Počítáme
s~posunem řádově 1\,cm, ověřené to není.

\begin{skyrentakeaway}
Hladké plochy dávají zrcadlový odraz bez posunu od stopy: chyba pod
centimetr, měřená je kolmá vzdálenost. Nejhorším cílem pro přesnost je
drsná nakloněná plocha: jedno měření kolísá o~15\,cm i~víc a~pomůže jen užší
svazek a~průměrování přes dráhu přiblížení.
\end{skyrentakeaway}

% ===========================================================================
\skyrenchapter[\skyrenmodeexplanation]{Teorie: anténa}{Anténa a čočka}
% ===========================================================================

\begin{skyrenglance}
Šířka svazku je daná průměrem apertury ve vlnových délkách, takže na 24\,GHz
se do $\varnothing$\,35\,mm vejde jeden svazek kolem 22$^\circ$. Plastová
hyperbolická čočka ho vytvoří z~malé antény levně. Zisk čočky ale omezuje
vysílací výkon: vyzářený výkon v~pásmu 24\,GHz nesmí přesáhnout 100\,mW.
\end{skyrenglance}

\section{Apertura, svazek a~zisk}
\skyrenlead{Tři vztahy, které určují velikost senzoru.}

Pro kruhovou aperturu o~průměru $D$ platí přibližně
\[
\theta_3 \approx (58\ \text{až}\ 70)\,\frac{\lambda}{D}, \qquad
G \approx \eta\left(\frac{\pi D}{\lambda}\right)^2, \qquad
R_\mathrm{ff} = \frac{2D^2}{\lambda}
\]
Šířka svazku $\theta_3$ ve stupních platí s~koeficientem 58 pro rovnoměrné
osvětlení, se 70 pro osvětlení klesající k~okraji. Činitel využití apertury
$\eta$ bývá 0,5--0,6. Vzdálenost $R_\mathrm{ff}$ je začátek dalekého pole.
Pro $\varnothing$\,35\,mm na 24\,GHz vychází $\theta_3 \approx 21$--23$^\circ$,
$G \approx 16$\,dBi a~$R_\mathrm{ff} \approx 20$\,cm.

\begin{skyrentable}{4}{Průměr apertury podle pásma a~požadované šířky svazku.}
\skyrenth{Pásmo} & \skyrenth{$\lambda$} & \skyrenth{$\varnothing$ pro 22$^\circ$} & \skyrenth{$\varnothing$ pro 10$^\circ$} \\
\midrule
24\,GHz & 12,4\,mm & $\approx$\,30--35\,mm & $\approx$\,75\,mm \\
60\,GHz & 5,0\,mm & $\approx$\,10--12\,mm & $\approx$\,30--35\,mm \\
\end{skyrentable}

Do $\varnothing$\,35\,mm se tedy na 24\,GHz vejde právě jeden svazek kolem
22$^\circ$. Proto návrh používá \textbf{jednu anténu} pro vysílání i~příjem.
Dvě antény vedle sebe by měly svazek zhruba 23$\times$42$^\circ$ a~podle
tabulky posunů v~dílu Odraz od plošného cíle několikanásobně větší posun i~kolísání. Svazek
10$^\circ$, který by pomohl na drsných nakloněných plochách, se do
$\varnothing$\,35\,mm vejde jen na 60\,GHz.

\section{Hyperbolická čočka}
\skyrenlead{Tvar plochy plyne z~Fermatova principu: všechny paprsky z~ohniska
musí mít za čočkou stejnou optickou dráhu.}

Zdroj leží v~ohnisku ve vzdálenosti $f$ před vrcholem čočky. Paprsek, který
dopadne na plochu v~bodě $(x, r)$, má urazit stejnou optickou dráhu jako
paprsek v~ose, aby za rovnou plochou vyšla rovinná vlna. Uvnitř materiálu
s~indexem lomu $n$ se dráha násobí $n$, takže
\[
\sqrt{x^2 + r^2} = f + n\,(x - f)
\]
Po dosazení $x = f + z$, kde $z$ je výška profilu nad vrcholem, vyjde
\[
r^2 = 2\,(n-1)\,f\,z + (n^2-1)\,z^2
\]
tedy hyperbola. Parametrický model podle tohoto vztahu je
v~\skyrencmd{hw/lens/lens\_24ghz.scad}.

\begin{skyrenfigure}{Hyperbolická čočka z~PETG ($n = 1{,}62$, $f = 28$\,mm,
$\varnothing$\,35\,mm), zvětšeno. Paprsky z~ohniska se na hyperbolické
ploše lámou tak, že za čočkou běží rovnoběžně.}{fig:cocka}
\begin{tikzpicture}[x=1.7mm, y=1.7mm, line width=0.8pt,
  every node/.style={font=\fontsize{8}{10}\selectfont}]
  \fill[skyrenNeutral400] (34.71,17.50) -- (33.21,15.00) -- (31.82,12.50)
    -- (30.57,10.00) -- (29.51,7.50) -- (28.70,5.00) -- (28.18,2.50)
    -- (28.00,0.00) -- (28.18,-2.50) -- (28.70,-5.00) -- (29.51,-7.50)
    -- (30.57,-10.00) -- (31.82,-12.50) -- (33.21,-15.00) -- (34.71,-17.50)
    -- (36.21,-17.50) -- (36.21,17.50) -- cycle;
  \draw[skyrenInk] (34.71,17.50) -- (33.21,15.00) -- (31.82,12.50)
    -- (30.57,10.00) -- (29.51,7.50) -- (28.70,5.00) -- (28.18,2.50)
    -- (28.00,0.00) -- (28.18,-2.50) -- (28.70,-5.00) -- (29.51,-7.50)
    -- (30.57,-10.00) -- (31.82,-12.50) -- (33.21,-15.00) -- (34.71,-17.50)
    -- (36.21,-17.50) -- (36.21,17.50) -- cycle;
  \foreach \r/\x in {0/28.00, 6/28.99, 12/31.56, 16/33.80, -6/28.99, -12/31.56, -16/33.80} {
    \draw[skyrenRed, line width=0.6pt] (0,0) -- (\x,\r) -- (48,\r);
  }
  \fill[skyrenInk] (0,0) circle (0.7);
  \node[anchor=east] at (-1,0) {ohnisko};
  \draw[skyrenNeutral600, line width=0.5pt] (0,-21) -- (28,-21);
  \draw[skyrenNeutral600, line width=0.5pt] (0,-20) -- (0,-22);
  \draw[skyrenNeutral600, line width=0.5pt] (28,-20) -- (28,-22);
  \node[anchor=north] at (14,-21.5) {$f = 28$\,mm};
  \draw[skyrenNeutral600, line width=0.5pt] (51,-17.5) -- (51,17.5);
  \draw[skyrenNeutral600, line width=0.5pt] (50,-17.5) -- (52,-17.5);
  \draw[skyrenNeutral600, line width=0.5pt] (50,17.5) -- (52,17.5);
  \node[anchor=west] at (52.5,0) {$D = 35$\,mm};
  \node[anchor=east] at (33.5,-19.5) {hyperbolická plocha};
  \node[anchor=west] at (36.8,-19.5) {rovná plocha};
\end{tikzpicture}
\end{skyrenfigure}

Obr.~\ref{fig:cocka} ukazuje hlavní vlastnost: hyperbolická plocha míří
k~anténě, rovná plocha k~cíli. Fázový střed antény musí ležet přesně
v~ohnisku. Posun o~pár milimetrů podél osy svazek rozostří, posun napříč ho
vychýlí.

\section{Ztráty, odrazy a~protiodrazná vrstva}
\skyrenlead{Tištěný plast má na 24\,GHz zanedbatelné ztráty. Víc vadí odrazy
na plochách, které se vracejí do přijímače.}

Na každé ploše se odrazí podíl výkonu
\[
|\Gamma|^2 = \left(\frac{n-1}{n+1}\right)^2
\]
pro PETG 5,6\,\% ($-12{,}5$\,dB), což dělá 0,25\,dB ztráty na plochu.
Odražená vlna se vrací do antény a~přidá se k~vazbě vysílače do přijímače.
Dielektrické ztráty jsou
\[
\alpha_\mathrm{dB} \approx 27{,}3\,\frac{n\tan\delta}{\lambda_0}
\]
na jednotku délky; pro PETG s~$\tan\delta \approx 0{,}01$ (odhad) to je
0,036\,dB/mm, ve středu čočky tlusté 8\,mm tedy kolem 0,3\,dB.

Odrazy odstraní \skyrenterm{protiodrazná vrstva} tloušťky $\lambda_0/4n_\mathrm{m}$
s~indexem lomu $n_\mathrm{m} = \sqrt{n}$. Pro PETG vychází $n_\mathrm{m} = 1{,}27$
($\varepsilon_\mathrm{r} = 1{,}62$) a~tloušťka 2,4\,mm. Tisk nabízí levný
trik: vrstvu vytisknout se sníženou výplní. Při lineárním směšování
permitivit s~vzduchem odpovídá $\varepsilon_\mathrm{r} = 1{,}62$ výplni kolem
38\,\%. Je to orientační výpočet; skutečnou permitivitu je třeba změřit.

\section{Buzení mimo osu a~přesah}
\skyrenlead{Anténa mimo ohnisko vychýlí svazek; záření, které čočku mine,
svítí do stran.}

Posune-li se fázový střed antény o~$d$ napříč osou, svazek se vychýlí zhruba
o~úhel $d/f$ na opačnou stranu. U~modulu s~oddělenou vysílací a~přijímací
anténou se tak jejich svazky pod jednou čočkou rozejdou. Při rozestupu
10\,mm a~$f = 28$\,mm to je $\pm$\,10$^\circ$: společný svazek se zúží
a~ztratí několik dB. Na vlastní desce s~jednou anténou tento problém odpadá.

Okraj čočky vidí anténa pod úhlem $\arctan(D/2f) = 32^\circ$. Co vyzáří pod
větším úhlem, čočku mine (\skyrenterm{přesah}) a~svítí do stran, kde vrací
odrazy od bližších ploch. Mezi anténu a~čočku proto patří trubka, zevnitř
vodivá nebo s~absorbérem.

\section{Vyzářený výkon}
\skyrenlead{Zisk čočky se počítá do vyzářeného výkonu. S~16\,dBi smí vysílač
dávat nejvýš 4\,dBm.}

V~pásmu 24,00--24,25\,GHz je pro zařízení krátkého dosahu povolený
ekvivalentní izotropně vyzářený výkon (\skyrenterm{EIRP}) 100\,mW, tedy
20\,dBm (ERC/REC 70-03, příloha~1; aktuální znění je třeba ověřit). EIRP je
součet vysílacího výkonu a~zisku antény, takže s~čočkou 16\,dBi smí vysílač
dávat nejvýš 4\,dBm. Hotový modul, který je certifikovaný s~vlastní anténou,
může po přidání čočky limit překročit. Vysílací výkon je pak nutné ve
firmwaru snížit. Rozpočet výkonu v~dílu Návrh už s~4\,dBm počítá.

\begin{skyrentakeaway}
Čočka $\varnothing$\,35\,mm dá na 24\,GHz svazek kolem 22$^\circ$ a~zisk
16\,dBi za cenu plastu. Její tvar plyne přímo z~Fermatova principu, ztráty
jsou zanedbatelné a~odrazy lze potlačit vrstvou s~řidší výplní. Zisk ale
snižuje povolený vysílací výkon na 4\,dBm.
\end{skyrentakeaway}

% ===========================================================================
\skyrenchapter[\skyrenmodereference]{Návrh}{Doporučené řešení}
% ===========================================================================

\begin{skyrenglance}
Pásmo 24\,GHz, jedna anténa s~čočkou $\varnothing$\,35\,mm, integrovaný
radarový čip s~trojúhelníkovou rampou a~levný mikrokontrolér. Tabulky
uvádějí profil vysílání, rozpočet výkonu, rozpočet chyby, čočku a~cenu.
Hodnoty čipu, které nejsou ověřené, jsou takto označené.
\end{skyrenglance}

\section{Blokové schéma}
\skyrenlead{Čtyři bloky: čočka s~anténou, vazebník, radarový čip
a~mikrokontrolér.}

\begin{skyrenfigure}{Blokové schéma senzoru. Vazebník spojuje vysílač
a~přijímač čipu do jedné antény; čip vzorkuje mezifrekvenční signál
a~mikrokontrolér z~něj počítá vzdálenost.}{fig:blok}
\begin{tikzpicture}[x=1mm, y=1mm, line width=0.8pt,
  every node/.style={font=\fontsize{7.5}{9}\selectfont, align=center}]
  \draw[skyrenInk] (0,0) rectangle (24,18);
  \node at (12,9) {čočka\\+ anténa\\$\varnothing$\,35\,mm};
  \draw[skyrenNeutral600, <->] (24,9) -- (31,9);
  \draw[skyrenInk] (31,0) rectangle (51,18);
  \node at (41,9) {vazebník\\(rat-race)};
  \draw[skyrenNeutral600, <-] (51,13) -- (59,13);
  \node[anchor=south] at (55,13) {Tx};
  \draw[skyrenNeutral600, ->] (51,5) -- (59,5);
  \node[anchor=north] at (55,5) {Rx};
  \draw[skyrenRed] (59,0) rectangle (93,18);
  \node at (76,9) {S3KM111L\\PLL, rampa $\uparrow\downarrow$,\\Tx, Rx, ADC};
  \draw[skyrenNeutral600, ->] (93,9) -- (100,9);
  \node[anchor=south] at (96.5,9) {SPI};
  \draw[skyrenInk] (100,0) rectangle (132,18);
  \node at (116,9) {MCU\\okno, FFT, interpolace,\\$\uparrow\downarrow$, tracking};
  \draw[skyrenNeutral600, ->] (132,9) -- (142,9);
  \node[anchor=south] at (137.5,9.5) {UART};
\end{tikzpicture}
\end{skyrenfigure}

Obr.~\ref{fig:blok} ukazuje cílovou vlastní desku. Při vývoji na hacknutém
modulu LD2450 místo vazebníku a~jedné antény slouží jeho oddělené antény
a~čočka sedí nad oběma.

\section{Profil vysílání}
\skyrenlead{Cílový profil pro dosah do 5\,m.}

\begin{skyrentable}{3}{Cílový profil. Zda ho integrovaný čip umí nastavit,
zejména délku chirpu, zatím nevíme.}
\skyrenth{Parametr} & \skyrenth{Hodnota} & \skyrenth{Důvod} \\
\midrule
Šířka přeladění & 250\,MHz & celé volné pásmo 24\,GHz \\
Rampa & trojúhelníková & vyruší Dopplerův posuv \\
Délka chirpu & 10\,\textmu s & $f_\mathrm{b}$ nad pásmem blikavého šumu \\
$f_\mathrm{b}$ na 1 / 5 / 8\,m & 167\,kHz / 833\,kHz / 1,33\,MHz & \\
Vzorkování & 4\,MSPS, I/Q & 40 vzorků, Hann, FFT 64 nebo 128 \\
Vysílací výkon & $\leq$\,4\,dBm & EIRP 20\,dBm se ziskem 16\,dBi \\
Výstup & 1 měření na pár chirpů & 4--10\,kHz podle čipu a~MCU \\
\end{skyrentable}

Integrovaný čip \skyrencmd{ICLegend S3KM111L} podle databriefu umí
trojúhelníkovou rampu a~až 8\,000 chirpů za sekundu, tedy kolem 4\,kHz měření.
Pokud ale nepodporuje chirp výrazně kratší než periodu 125\,\textmu s, klesne
rozdílová frekvence na 1\,m na desítky kHz a~měření pod 1--2\,m se zhorší.
Na 5\,m to nevadí.

\section{Rozpočet výkonu}
\skyrenlead{Na 5\,m je signálu dost i~pro nejslabší povrch.}

\begin{skyrentable}{3}{Rozpočet výkonu na 5\,m. Vysílač 4\,dBm, čočka
16\,dBi, vazebník 3,5\,dB v~každém směru, šum $-107$\,dBm na buňku (šumové
číslo 15\,dB je odhad).}
\skyrenth{Cíl} & \skyrenth{Přijatý výkon} & \skyrenth{SNR po FFT} \\
\midrule
plech, zrcadlově & $-51$\,dBm & 56\,dB \\
hladký beton, zrcadlově & $-59$\,dBm & 48\,dB \\
sádrokarton, zrcadlově & $-63$\,dBm & 44\,dB \\
suché dřevo, zrcadlově & $-66$\,dBm & 41\,dB \\
drsný povrch, $\sigma^0 = -10$\,dB & $-77$\,dBm & 30\,dB \\
drsný povrch, $\sigma^0 = -20$\,dB & $-87$\,dBm & 20\,dB \\
\end{skyrentable}

Mezi nejsilnějším případem (plech zblízka) a~nejslabším (drsný povrch
na 5\,m) je přes 60\,dB. K~tomu vazba vysílače do přijímače kolem
$-15$\,dBm. Proto horní propust s~+20\,dB na dekádu a~podle potřeby řízení
zesílení.

\section{Rozpočet chyby}
\skyrenlead{Na hladkých plochách pod centimetr, na drsných rozhoduje speckle
a~průměrování přes dráhu přiblížení.}

\begin{skyrentable}{4}{Rozpočet chyby na 5\,m se svazkem 22$^\circ$. Hodnoty
„po průměrování“ platí pro průměr přes 1,2\,m přiblížení (5\,ms při
240\,m/s), šum pro drsný povrch s~$\sigma^0 = -15$\,dB. Jde o~výpočty
a~simulace, ne o~měření.}
\skyrenth{Zdroj chyby} & \skyrenth{Hladký} & \skyrenth{Drsný, kolmo} & \skyrenth{Drsný, 15$^\circ$} \\
\midrule
šum, jedno měření & 0,1--0,4\,cm & 1,9\,cm & 1,9\,cm \\
posun od stopy po kalibraci & 0 & $\approx$\,1\,cm & $\approx$\,1--1,5\,cm \\
speckle, jedno měření & 0 & 7\,cm & 16\,cm \\
Dopplerův posuv po průměru ramp & $\approx$\,0 & $\approx$\,0 & $\approx$\,0 \\
\midrule
\textbf{celkem, jedno měření} & $<$\,0,5\,cm & \textbf{$\approx$\,7\,cm} & \textbf{$\approx$\,16\,cm} \\
\textbf{celkem po průměrování} & $<$\,0,5\,cm & \textbf{$\approx$\,3\,cm} & \textbf{$\approx$\,5--6\,cm} \\
\end{skyrentable}

Zadání 5\,cm tedy splní hladké plochy vždy a~drsné plochy kolmo k~ose po
průměrování. Drsná plocha nakloněná o~15$^\circ$ vychází na hranici. Pomohl
by užší svazek (s~10$^\circ$ vychází model na 3\,cm), ten se ale do
$\varnothing$\,35\,mm na 24\,GHz nevejde. Druhou možností je delší dráha
průměrování, pokud to aplikace dovolí.

\section{Čočka}
\skyrenlead{Parametry pro tisk a~obrábění.}

\begin{skyrentable}{4}{Čočka $\varnothing$\,35\,mm, ohnisková vzdálenost
28\,mm, okraj 1,5\,mm. Indexy lomu tištěných plastů jsou orientační.}
\skyrenth{Materiál} & \skyrenth{$n$} & \skyrenth{Tloušťka ve středu} & \skyrenth{Výška od antény} \\
\midrule
HDPE & 1,52 & 9,3\,mm & 37\,mm \\
PTFE & 1,45 & 10,3\,mm & 38\,mm \\
PETG (tisk, 100\,\% výplň) & $\approx$\,1,62 & 8,2\,mm & 36\,mm \\
PLA (tisk, 100\,\% výplň) & $\approx$\,1,65 & 7,9\,mm & 36\,mm \\
\end{skyrentable}

Očekávaná šířka svazku je 21--23$^\circ$, zisk kolem 16\,dBi a~daleké pole
začíná zhruba ve 20\,cm. Pro první prototypy stačí tisk z~PETG; pro další
iterace je lepší HDPE se známými vlastnostmi a~nulovými ztrátami.

\section{Cena}
\skyrenlead{Materiál na kus při 500 kusech. Rozhoduje radarový čip.}

\begin{skyrentable}{3}{Varianty materiálu na kus. Cena čipu S3KM111L je
odhad z~cen hotových modulů, ne nabídka.}
\skyrenth{Varianta} & \skyrenth{Radarový čip} & \skyrenth{Celkem} \\
\midrule
Integrovaný SoC (doporučeno) & S3KM111L, $\approx$\,1--2\,€ (neověřeno) & $\approx$\,4--7\,€ \\
Záloha s~vlastním front-endem & BGT24MTR11, $\approx$\,10\,€ & $\approx$\,15\,€ \\
60\,GHz s~anténou v~pouzdře & Acconeer A121, $\approx$\,9\,\$ & $\approx$\,10--12\,€ \\
\end{skyrentable}

Integrovaný čip obsahuje vysílač, přijímač, PLL, generátor rampy a~ADC,
takže odpadá rychlý ADC, generátor rampy i~teplotně kompenzovaný krystal.
Rampa z~PLL je lineární, takže odpadá i~kalibrace linearity při výrobě.
Stačí levný mikrokontrolér, který přijímá data a~počítá malé FFT. Anténa
může být na FR4 nebo levném PTFE laminátu F4B: ztráty 2--3\,dB a~rozladění
kvůli toleranci permitivity jsou při rezervě přes 20\,dB přijatelné.

\begin{skyrenfigure}{Volba varianty. O~velikosti rozhoduje, kterou
vzdálenost aplikace potřebuje.}{fig:volba}
\begin{tikzpicture}[x=1mm, y=1mm, line width=1pt,
  every node/.style={font=\fontsize{8}{10}\selectfont}]
  \draw[skyrenInk] (-24,0) rectangle (24,10);
  \node at (0,5) {Kterou vzdálenost potřebuješ?};
  \draw[skyrenNeutral600, line width=0.6pt] (0,0) -- (0,-6);
  \draw[skyrenNeutral600, line width=0.6pt] (-36,-6) -- (36,-6);
  \draw[skyrenNeutral600, line width=0.6pt] (-36,-6) -- (-36,-13);
  \draw[skyrenNeutral600, line width=0.6pt] (36,-6) -- (36,-13);
  \node[anchor=south] at (-36,-5) {v~ose senzoru};
  \node[anchor=south] at (36,-5) {k~nejbližšímu povrchu};
  \draw[skyrenRed] (-58,-27) rectangle (-14,-13);
  \node at (-36,-17) {úzký svazek $\approx$\,22$^\circ$};
  \node at (-36,-22.5) {24\,GHz + čočka $\varnothing$\,35, 4--7\,€};
  \draw[skyrenInk] (14,-27) rectangle (58,-13);
  \node at (36,-17) {široký svazek stačí};
  \node at (36,-22.5) {A121 bez čočky, $\approx$\,8$\times$8\,mm};
\end{tikzpicture}
\end{skyrenfigure}

Široký svazek přirozeně měří nejbližší bod plochy. Pokud aplikaci stačí
nejkratší vzdálenost k~povrchu, úzký svazek není potřeba a~velikost určí
jen čip (obr.~\ref{fig:volba}). Pro nejbližší bod je ale potřeba jemné
rozlišení, aby šla přesně najít náběžná hrana odrazu --- proto v~pravé
větvi 60\,GHz, ne 24\,GHz s~buňkou 60\,cm.

\begin{skyrentakeaway}
Doporučená varianta stojí na integrovaném 24GHz čipu za jednotky eur, jedné
anténě přes $\varnothing$\,35\,mm a~tištěné čočce. Na hladkých plochách měří
pod centimetr, na drsných kolmých po průměrování kolem 3\,cm. Drsné nakloněné
plochy a~rozšíření na 3\,cm do 30\,m potřebují užší svazek, v~tomto průměru
tedy 60\,GHz.
\end{skyrentakeaway}

% ===========================================================================
\skyrenchapter[\skyrenmodehowto]{Postup}{Čip z~levného modulu}
% ===========================================================================

\begin{skyrenglance}
Dokumentace k~čipu S3KM111L není veřejná. Rychlejší cesta než čekat na NDA
je rozebrat hotový modul HLK-LD2450, vytáhnout z~jeho mikrokontroléru
inicializaci čipu a~nahradit firmware vlastním. Modul slouží jako vývojová
platforma; finální tvar $\varnothing$\,35\,mm bude vlastní deska.
\end{skyrenglance}

\section{Co koupit}
\skyrenlead{Dva moduly z~lokální nabídky a~vývojové pomůcky.}

\begin{skyrentable}{3}{Nákupní seznam. Moduly jsou skladem v~ČR.}
\skyrenth{Položka} & \skyrenth{Počet} & \skyrenth{Účel} \\
\midrule
HLK-LD2450 (FMCW, 1T2R, 44$\times$15\,mm) & 2--3 & hlavní cíl rozboru \\
HLK-LD2451 (FMCW, $\pm$\,20$^\circ$, 70$\times$35\,mm) & 1 & reference laděná na rychlé cíle \\
Převodník USB--UART & 1 & čtení výstupu modulu \\
Ladicí sonda SWD (ST-Link, Pico) & 1 & čtení a~zápis firmwaru \\
Logický analyzátor & 1 & odposlech rozhraní čip--MCU \\
\end{skyrentable}

Moduly \skyrencmd{HLK-LD2415H}, \skyrencmd{HW-MS03} a~\skyrencmd{RCWL-0516}
nepoužíváme. Jsou to CW dopplerovské senzory, měří jen pohyb nebo rychlost,
vzdálenost z~nich žádné zpracování nedostane.

\section{Kroky}
\skyrenlead{Pořadí je závazné --- každý krok odstraňuje nejistotu, na které
stojí další.}

\begin{enumerate}
\item \textbf{Nafotit obě strany LD2450} a~určit radarový čip, mikrokontrolér
      a~piny SWD. Z~fotek se zároveň určí poloha vysílací a~přijímací antény
      pro umístění čočky.
\item \textbf{Vytáhnout firmware přes SWD.} Obsahuje inicializační sekvenci
      čipu, tedy registry, ke kterým chybí dokumentace. Pokud je paměť
      zamčená, zbývá odposlech rozhraní.
\item \textbf{Odposlechnout rozhraní mezi čipem a~MCU} logickým analyzátorem
      při startu a~během měření. Cílem je zjistit, jak se nastavuje rampa,
      vysílací výkon a~v~jakém formátu čip posílá vzorky.
\item \textbf{Vytisknout čočku} z~PETG podle \skyrencmd{lens\_24ghz.scad}
      a~změřit index lomu (viz níže). Podle něj model přepočítat.
\item \textbf{Nahrát vlastní firmware}: trojúhelníková rampa přes 250\,MHz,
      vysílací výkon snížený na limit EIRP, surová data, Hannovo okno, FFT,
      parabolická interpolace, kombinace ramp, vyřazení slabých měření
      a~tracking.
\item \textbf{Změřit} šířku svazku s~čočkou a~přesnost na plechu, hladké
      desce a~cihlové zdi, vždy kolmo a~při náklonu 15$^\circ$. Cihla ověří
      model speckle. Teprve tato čísla nahradí výpočty v~dílu Návrh.
\item \textbf{Navrhnout vlastní desku $\varnothing$\,35\,mm} se stejným čipem
      a~jednou anténou pod čočkou.
\end{enumerate}

\begin{skyrenexample}
Index lomu změří radar sám. Kovová deska stojí ve známé vzdálenosti; do
svazku se vloží rovná destička z~materiálu čočky o~tloušťce 10\,mm.
Zdánlivá vzdálenost se prodlouží o~$(n-1)\cdot t$. Prodloužení 6,2\,mm
znamená $n = 1{,}62$.
\end{skyrenexample}

\section{Na co dát pozor}
\skyrenlead{Čtyři věci, které rozhodnou, jestli prototyp s~čočkou funguje.}

\textbf{Rozestup antén modulu.} Vysílací a~přijímací anténa LD2450 jsou vedle
sebe, takže se jejich svazky pod jednou čočkou rozejdou (díl Anténa
a~čočka). Osu čočky
je třeba umístit doprostřed mezi vysílač a~bližší přijímač. Rozestup přes
20\,mm by byl problém.

\textbf{Přesah.} Mezi modul a~čočku patří trubka, zevnitř buď s~hliníkovou
páskou, nebo s~absorbérem --- která varianta je lepší, rozhodne měření.

\textbf{Vyzářený výkon.} S~čočkou 16\,dBi smí vysílač dávat nejvýš 4\,dBm.
Pokud firmware modulu vysílá víc, je nutné výkon snížit.

\textbf{Shoda CE/RED.} Hotový modul má shodu pro svůj původní firmware
a~anténu. Změnou profilu vysílání nebo přidáním čočky přechází odpovědnost
za shodu na nás. Pro vývoj to nevadí, pro prodej ano.

\begin{skyrentakeaway}
Rozbor modulu LD2450 je nejlevnější cesta k~dokumentaci čipu i~k~prvnímu
změřenému prototypu. Výsledkem prvního kola mají být čísla, která nahradí
výpočty a~simulace v~tomto dokumentu --- hlavně chování na drsné nakloněné
ploše --- a~podklady pro vlastní desku $\varnothing$\,35\,mm.
\end{skyrentakeaway}

% ===========================================================================
\skyrenappendix{Přehled vzorců}
% ===========================================================================

Vztahy z~teoretických dílů a~z~návrhu na jednom místě. $c$ je rychlost světla,
$\lambda$ vlnová délka, $B$ šířka přeladění, $T$ délka chirpu, $R$ vzdálenost.

\begin{skyrendeflist}{Veličina}{Vztah}
rozdílová frekvence & $f_\mathrm{b} = 2RB/(cT)$ \\
rozlišení & $\Delta R = c/2B$ \\
Cramér--Raova mez & $\sigma_R \geq 0{,}39\,\Delta R/\sqrt{\mathrm{SNR}}$; prakticky $0{,}6\,\Delta R/\sqrt{\mathrm{SNR}}$ \\
parabolická interpolace & $\delta = \tfrac{1}{2}(\alpha-\gamma)/(\alpha-2\beta+\gamma)$ \\
dopplerovský posuv & $f_\mathrm{d} = 2v/\lambda$ \\
Dopplerův posuv odhadu & $\Delta R_\mathrm{D} = vcT/(\lambda B)$ \\
kombinace ramp & $R = (R_\uparrow+R_\downarrow)/2$, $v = \lambda B (R_\downarrow-R_\uparrow)/(2cT)$ \\
šum na buňku & $N_\mathrm{dBm} = -174 + 10\log_{10}(1{,}5/T) + F_\mathrm{dB}$ \\
bodový cíl & $P_\mathrm{r} = P_\mathrm{t}G^2\lambda^2\sigma/((4\pi)^3R^4)$ \\
zrcadlová plocha & $P_\mathrm{r} = P_\mathrm{t}G^2\lambda^2|\Gamma|^2/((4\pi)^2(2R)^2)$ \\
drsná plocha & $\sigma = \sigma^0 R^2\,\pi\theta_3^2/(8\ln 2)$, $P_\mathrm{r} \propto G/R^2$ \\
činitel odrazu kolmo & $\Gamma = (1-\sqrt{\varepsilon_\mathrm{r}})/(1+\sqrt{\varepsilon_\mathrm{r}})$ \\
Rayleighovo kritérium & $\rho = \exp[-(4\pi h\cos\theta_\mathrm{i}/\lambda)^2]$ \\
první Fresnelova zóna & $r_1 = \sqrt{\lambda R/2}$ \\
posun od stopy svazku & $b = R\theta_3^2(1+\tan^2\alpha)/(16\ln 2)$ \\
kolísání (speckle) & $s_R \approx R\sigma_\theta\sqrt{\sigma_\theta^2+\tan^2\alpha}$, $\sigma_\theta = \theta_3/\sqrt{16\ln 2}$ \\
průměr přes dráhu $L$ & $\sigma_\mathrm{prům} \approx \sqrt{s_R\lambda R/(4L)}$ \\
šířka svazku & $\theta_3 \approx (58\ \text{až}\ 70)\,\lambda/D$ [$^\circ$] \\
zisk apertury & $G \approx \eta(\pi D/\lambda)^2$ \\
daleké pole & $R_\mathrm{ff} = 2D^2/\lambda$ \\
profil čočky & $r^2 = 2(n-1)fz + (n^2-1)z^2$ \\
odraz na ploše čočky & $|\Gamma|^2 = ((n-1)/(n+1))^2$ \\
protiodrazná vrstva & $n_\mathrm{m} = \sqrt{n}$, tloušťka $\lambda_0/(4n_\mathrm{m})$ \\
EIRP & $P_\mathrm{t} + G \leq 20$\,dBm (24,00--24,25\,GHz) \\
\end{skyrendeflist}

% ===========================================================================
\skyrenappendix{Zvažované a zamítnuté varianty}
% ===========================================================================

Během studie jsme prošli řadu cest. Tabulka shrnuje, proč nevyhrály,
aby se k~nim nebylo třeba vracet.

\begin{skyrendeflist}{Varianta}{Proč nevyhrála}
Parkside PLEM 20 A4 & fázový měřák s~uzavřeným ASICem, frekvence i~integrace
  zadrátované v~čipu \\
optický dToF (LightWare SF30/D) & zadání splní, cenou patří mezi měřicí
  techniku \\
levné dToF senzory (VL53Lxx, TMF8828, TF-Luna) & desítky až stovky Hz,
  vícezónové čipy navíc integrují desítky ms \\
hotové radarové moduly (NanoRadar NRA24, MR72) & neznámé chování firmwaru
  při 240\,m/s, cena a~velikost \\
vyřazené automobilové radary (ARS408, ESR) & 77\,GHz, filtry věrohodnosti
  ve firmwaru, velikost \\
24\,GHz s~polem 8$\times$8 & velikost 130$\times$70\,mm, do $\varnothing$\,35
  se nevejde \\
diskrétní front-end na 60/77\,GHz & jednotlivé obvody stojí desítky eur,
  vyžadují speciální laminát a~bez měřicí techniky do 80\,GHz je nelze
  ověřit \\
60\,GHz s~anténou v~pouzdře (IWRL6432AOP) & technicky vhodný i~pro 30\,m
  s~čočkou $\varnothing$\,35, cenu se nepodařilo zjistit \\
Acconeer A121 & nejmenší varianta, ale zhruba 3$\times$ dražší; chování při
  240\,m/s neověřené, protože vzdálenostní body snímá postupně \\
radar se stojatým vlněním & stejné meze jako FMCW; výhodou je jediná anténa,
  nevýhodou malý signál na velké stejnosměrné složce a~vícenásobné odrazy \\
FSK (několik diskrétních tónů) & milimetrová přesnost pro jeden cíl, ale
  měří jen za pohybu \\
UWB čipy jako radar (DW3000) & velké antény a~široký svazek na 6,5--8\,GHz \\
využít 1\,GHz modulace S3KM111L & mimo pásmo 24,00--24,25\,GHz je vysílání
  v~EU nelegální \\
\end{skyrendeflist}

% ===========================================================================
\skyrenappendix{Zdroje a vstupní data}
% ===========================================================================

\subsection{Co nebylo ověřeno}
Dokument stojí na katalogových údajích výrobců, veřejné dokumentaci,
výpočtech a~simulacích. \textbf{Nic z~uvedených dosahů, přesností, šířek
svazku ani cen nebylo ověřeno vlastním měřením ani poptávkou.} Konkrétně
chybí:

\begin{itemize}
\item parametry čipu S3KM111L --- minimální délka chirpu, šířka pásma IF,
      vysílací výkon, šumové číslo, kompresní bod přijímače, cena
      a~dostupnost dokumentace,
\item koeficienty zpětného rozptylu $\sigma^0$, drsnosti a~permitivity
      povrchů na 24\,GHz (orientační hodnoty),
\item index lomu a~ztrátový činitel tištěných plastů, poloha fázových středů
      antén LD2450,
\item model speckle je ověřený jen simulací, ne měřením na skutečné zdi,
\item chování A121 při rychlosti 240\,m/s,
\item regulatorní rámec: limit EIRP podle ERC/REC 70-03 a~harmonizované
      normy (ETSI EN 305~550 a~související) v~aktuálním znění.
\end{itemize}

\subsection{Vstupní data a~simulace}
Fotodokumentace teardownu Parkside PLEM 20 A4 je ve složce \skyrencmd{img/},
model čočky v~\skyrencmd{hw/lens/}. Simulace interpolace vrcholu a~speckle
na drsné zdi jsou ve složce \skyrencmd{sim/}; jsou napsané v~čistém Pythonu
bez dalších knihoven.

\subsection{Zdroje}
\begin{skyrendeflist}{Zdroj}{Co z~něj pochází}
ICLegend, S3KM111L databrief (Sekorm) & integrovaný 24GHz čip:
  trojúhelníková rampa, až 8\,kHz chirpů, QFN-32 4$\times$4\,mm \\
Hi-Link, HLK-LD2450 a~LD2451 & parametry modulů: FMCW, 1T2R, rozměry,
  šířka svazku a~přeladění LD2451 \\
Hi-Link, HLK-LD2415H & CW dopplerovský modul, důvod vyřazení \\
Infineon, BGT24MTR11 & 24GHz transceiver pro záložní variantu, ceny RS \\
Acconeer, A121 & 60GHz pulzní koherentní radar, ceny LCSC a~DigiKey \\
Infineon, BGT60TR13C & 60GHz čip s~anténou v~pouzdře, šířka pásma 5,5\,GHz \\
TI, IWRL6432AOP & 60GHz FMCW s~anténou v~pouzdře, IF 5\,MHz, ADC
  12,5\,MSPS \\
TI, AWR2944 a~IWR6843 & šířka přeladění 4--4,5\,GHz v~pásmech 77 a~60\,GHz \\
LightWare, SF30/D & optický dálkoměr 200\,m, až 20\,000 měření/s \\
Broadcom, AFBR-S50MV85G & optický dToF: 10\,m, až 3000 snímků/s \\
nanoradar-protocols, tier4/ars408\_driver (GitHub) & protokoly hotových
  a~automobilových radarů \\
ERC/REC 70-03 & limit EIRP 100\,mW pro 24,00--24,25\,GHz (neověřené
  aktuální znění) \\
\end{skyrendeflist}

% ===========================================================================
\skyrenappendix{Slovníček pojmů}
% ===========================================================================

\begin{skyrendeflist}{Pojem}{Význam jednou větou}
FMCW & radar se spojitou vlnou, jejíž kmitočet lineárně přelaďujeme;
  vzdálenost vypadne z~rozdílu kmitočtů \\
chirp & jedna kmitočtová rampa vysílaného signálu \\
trojúhelníková rampa & střídání stoupajícího a~klesajícího chirpu; ruší
  Dopplerův posuv odhadu vzdálenosti \\
rozdílová frekvence & rozdíl mezi vyslaným a~přijatým kmitočtem, přímo úměrný
  vzdálenosti cíle \\
I/Q vzorkování & vzorkování dvou kanálů posunutých o~90$^\circ$; dává
  komplexní signál se znaménkem kmitočtu \\
okno & váhová funkce vzorků před FFT, která potlačí postranní laloky spektra \\
parabolická interpolace & dopočet polohy vrcholu mezi body FFT z~tří hodnot \\
rozlišení & nejmenší vzdálenost, na které radar oddělí dva cíle; $c/2B$ \\
přesnost & jak přesně radar určí polohu jednoho cíle; omezuje ji odstup
  signálu od šumu \\
Cramér--Raova mez & teoretická nejmenší chyba odhadu při daném odstupu
  signálu od šumu \\
Dopplerův posuv & změna kmitočtu odrazu pohybujícího se cíle; u~FMCW posune
  odhad vzdálenosti \\
šumové číslo & o~kolik přijímač zhorší odstup signálu od šumu proti
  ideálnímu \\
blikavý šum (1/$f$) & šum směšovače, který roste k~nízkým kmitočtům \\
plošný cíl & cíl, který vyplní celý svazek; přijatý výkon klesá s~$R^2$ \\
zrcadlový odraz & koherentní odraz od hladké plochy, vzniká v~první
  Fresnelově zóně \\
difúzní rozptyl & odraz od drsné plochy do všech směrů \\
první Fresnelova zóna & skvrna kolem paty kolmice, ze které vzniká
  zrcadlový odraz \\
Rayleighovo kritérium & mez drsnosti, nad kterou se plocha přestává chovat
  jako zrcadlo \\
koeficient zpětného rozptylu & $\sigma^0$, radarový průřez drsné plochy na
  jednotku plochy \\
HPBW & šířka svazku v~bodech polovičního výkonu \\
stopa svazku & část plochy ozářená svazkem \\
posun od stopy svazku & systematický posun odhadu za kolmou vzdálenost
  u~drsné plochy \\
speckle & náhodné kolísání odhadu kvůli sčítání příspěvků drsné plochy
  s~náhodnými fázemi \\
hyperbolická čočka & čočka, jejíž plocha převede kulovou vlnu z~ohniska na
  rovinnou \\
protiodrazná vrstva & vrstva tloušťky čtvrt vlny, která potlačí odraz na
  ploše čočky \\
přesah & část vyzářeného výkonu, která čočku mine \\
EIRP & ekvivalentní izotropně vyzářený výkon; součet vysílacího výkonu
  a~zisku antény \\
dToF & přímé měření doby letu pulzu \\
fázový měřák & měření vzdálenosti z~fázového posuvu modulovaného svazku \\
pásmo ISM 24\,GHz & volně použitelné pásmo 24,00--24,25\,GHz, šířka 250\,MHz \\
\end{skyrendeflist}

\end{document}
